% !TeX root = paper.tex
\clearpage\onecolumn\section*{Version History}
\begingroup
\small
\renewcommand{\arraystretch}{1.2}
\begin{longtable}{|l|p{0.14\linewidth}|p{\dimexpr0.86\linewidth-4em-6\tabcolsep-4\arrayrulewidth\relax}|}
\hline
\textbf{Version} & \textbf{Date} & \textbf{Changes} \\ \hline
\endhead
0.1 & 7 October 2026 & Starting file: IEEE layout and section headings. \\ \hline
0.2 & 8 October 2026 & Sections III (Data) and IV (Labels and Splits) drafted. \\ \hline
0.3 & 9 October 2026 & Section V (Preprocessing) drafted \\ \hline
0.4 & 23 October 2026 & Sections VI (Model and Training), VII (Calibration, Thresholds and Abstention), IX (Application) and X (Limitations) drafted from the repository. Missing full stop added in Section V. \\ \hline
0.5 & 28 October 2026 & Sections VIII-A (Internal Test, with Table~\ref{tab:I}), VIII-C (Multi-Task vs Single-Task Ablation) and VIII-D (Calibration and Abstention) drafted from the Task 13 results. External protocol paragraph added to VIII-B. VIII-E renamed and media-opacity haze check added from PR 85. Two limitations added. \\ \hline
0.6 & 29 October 2026 & Sections I (Introduction) and II (Literature Review) drafted. Citations numbered in order of first use and the reference list added (28 entries, each checked against PubMed, Crossref or the publisher). \\ \hline
0.7 & 30 October 2026 & VIII-E: synthetic blur result added from PR 86 (open); media-opacity figures confirmed against merged PR 85. Haze limitation and one sentence in Section II updated to match. \\ \hline
0.8 & 31 October 2026 & VIII-B: external results and Table~\ref{tab:II} added from PR 87 (open). VIII-E: Grad-CAM heatmap paragraph added from PR 88 (open); blur figures confirmed against merged PR 86. Abstract drafted. Limitations updated for the external results, and one new limitation added. \\ \hline
0.9 & 3 November 2026 & VIII-E: final haze verdict added from SCRUM-100 (PR 89); heatmap paragraph confirmed against merged PR 88. All Task 15 results are now on main (PR 91). Two-column layout restored for the body (lost in v0.6). \\ \hline
0.10 & 6 November 2026 & VIII-E: optic disc heatmap check added from Task 16 (PR 92, merged to develop). VIII-B: external figures confirmed against merged PR 87. Section XI (Conclusion and Future Work) drafted. Abstract and one limitation updated. Section IX checked against the merged back end (PR 90). All 28 references rechecked against Crossref, PubMed, arXiv, PMLR and GitHub; titles corrected for \cite{ref28} (published chapter title) and wording for \cite{ref24}. Every cited claim checked against its source; wording tightened for \cite{ref6}, \cite{ref17} and \cite{ref18}, and the LNCS volume added to \cite{ref16}. Final clean-up: all open notes resolved or removed, Fig.~\ref{fig:1} (pipeline diagram) added, training details and the AI use disclosure written, and contribution (6) corrected to match Section IX. Review fixes: application described as a prototype with a mock model, freeze and run records cited by commit, protocol draft status and the earlier 21-eye trial disclosed, dataset exclusion list shortened, fellow-eye design explained, retina\_dataset\_2016 limits stated, optic disc baseline qualified, AI disclosure revised, and some wording simplified. \\ \hline
0.11 & 10 November 2026 & Task 17: abstract, VIII-C and conclusion updated to match the shared versus single-task comparison in PR 95 (open, 9 November); all its figures already matched Tables I and II. 9 November: every reported figure rechecked against the repository (main at 1222b7d); one sentence in VIII-B corrected, since the external differences between designs were not mostly within the spread between seeds; wording edited throughout for clarity and consistency. Layout aligned with the IEEE template: two empty paragraphs removed from Section IV, Section V spacing matched to the rest, table text, captions and references set to 8 pt, abstract text no longer italic. 10 November: points from the team review addressed. IV now says where the other two cataract eyes are. VII explains the uncertain label and the choice of seed 42. VIII-B corrects the ACRIMA wording, and VIII-E states the eligibility rules and both warning criteria. Section IX separates back-end and front-end features. Table~\ref{tab:II} columns widened (rebalanced again for 9 pt text). The abstract wording on cataract was revised, and pending notes in square brackets were added to the yellow parts. Contribution (5) now lists all shortcut checks. VIII-C now reports the paired AUROC differences from the Task 17 analysis (PR 95, head 7178102, approved 10 November). Formatting matched to the course IEEE template: paragraphs indented with no extra space between them, table text 9 pt, table captions on two lines in title case, journal and proceedings names in italic, author emails added in 9 pt Courier, title block spacing matched, and Times New Roman set as the font for the whole document. All spacing then checked element by element against the template styles: heading, caption, figure, abstract and reference spacing matched, level-2 headings and references use hanging indents, and table header cells are centred after the first column. Tables I and II are set to stay on one page with their captions. Two figures added to VIII-E from the merged Task 15 and Task 16 results: Fig.~\ref{fig:1}, the synthetic blur test, redrawn at column size from results/task15/synthetic\_blur\_summary.csv, and Fig.~\ref{fig:2}, a Grad-CAM panel built from results/task15/cataract\_heatmaps and results/task16/glaucoma\_heatmap\_contact\_sheet.png. The pipeline figure is now Fig.~\ref{fig:3}. Page size changed from US Letter to A4 to match the course template, with its margins (0.75 in top, 0.56 in sides, 1.69 in bottom) and column gap; Table~\ref{tab:II} and this table narrowed slightly to fit. Table~\ref{tab:II} moved below the last paragraph of VIII-B and Fig.~\ref{fig:2} moved above the optic disc paragraph, so that neither is split or leaves a gap in a column. \\ \hline
\end{longtable}
\endgroup
